\documentclass[../../../../main.tex]{subfiles}
\begin{document}

\begin{flushright}
\footnotesize\textit{【自動文字起こし・要確認】}
\end{flushright}

[解] $c = \cos\theta \ (-1 \le c \le 1)$ とおくと
\begin{equation*}
(\text{与式}) \iff a(2c^2-1) + bc - 1 < 0 \quad \cdots \text{\textcircled{1}}
\end{equation*}
\textcircled{1}の左辺$f(c)$とおき、$f(c)<0$ となる条件をしらべる。\\
$f'(c) = 2ac+b$ である。

$1^\circ \ a=0$の時、$f(c)$は高々一次関数だから
\begin{align*}
&f(1) < 0 \land f(-1) < 0 \\
\iff & -1 < b < 1
\end{align*}

$2^\circ \ a>0$の時\\
$f(c)$の軸 $c = -\frac{b}{4a}$ によらず、$f(c)$は $c = \pm 1$で最大で
\begin{align*}
&f(1) < 0 \land f(-1) < 0 \\
\iff & a+b-1 < 0 \land a-b-1 < 0
\end{align*}

$3^\circ \ a<0$の時
\begin{align*}
&\text{\textcircled{ア}} \ -1 \le -\frac{b}{4a} \le 1 \text{の時} \quad f(-\frac{b}{4a}) < 0 \iff \frac{(a+\frac{1}{2})^2}{1/4} + \frac{b^2}{1} < 1 \quad \cdots \text{A} \\
&\text{\textcircled{イ}} \ -\frac{b}{4a} \le -1 \text{の時} \quad f(-1) < 0 \\
&\text{\textcircled{ウ}} \ -\frac{b}{4a} \ge 1 \text{の時} \quad f(1) < 0
\end{align*}

以上まとめて
\begin{itemize}
\item $a \ge 0 \land a-1 < b < a+1$
\item $4a \le b \le -4a \land a < 0 \land \frac{(a+\frac{1}{2})^2}{1/4} + \frac{b^2}{2} < 1$
\item $4a \ge b \land a < 0 \land a-1 < b$
\item $-4a \le b \land a < 0 \land b < -a+1$
\end{itemize}

これを図示して下図(境界含まず)

\begin{center}
\begin{tikzpicture}[scale=1.5]
\draw[->] (-2,0) -- (1.5,0) node[right] {$a$};
\draw[->] (0,-2.5) -- (0,2.5) node[above] {$b$};
\draw[dashed] (-1, -2) -- (0.5, 2.5) node[right] {$b=4a$};
\draw[dashed] (-1, 2) -- (0.5, -2.5) node[right] {$b=-4a$};
\draw (-1.5, -0.5) -- (0,1) node[right] {$b=a+1$};
\draw (-1.5, 0.5) -- (0,-1) node[right] {$b=-a-1$};
\draw (0,1) -- (1,2) node[right] {};
\draw (0,-1) -- (1,-2) node[right] {};
\draw[domain=-1:0, smooth, variable=\a, blue] plot ({\a}, {sqrt(2*(1 - 4*(\a+0.5)^2))});
\draw[domain=-1:0, smooth, variable=\a, blue] plot ({\a}, {-sqrt(2*(1 - 4*(\a+0.5)^2))});
\node at (-0.5, 0) [below] {$-\frac{1}{2}$};
\node at (-1, 0) [below left] {$-1$};
\node at (0, 0) [below right] {$O$};
\node at (-0.33, 0) [below] {$-\frac{1}{3}$};
\node[blue] at (1.5, 1) {[本時のミス]};
\node[blue] at (1.5, 0.5) {$\cdot$ 計算ミス};
\node[blue] at (1.5, 0) {$\cdot$ 軸についての勘違い};
\node at (-1, 1.5) {$\frac{(a+1/2)^2}{1/4} + \frac{b^2}{2} = 1$};
\node at (2.5, 2) {$4(a+1/2)^2 + \frac{b^2}{2} = 1, \ b = \pm (a-1) \text{ を連立すると}$};
\node at (2.5, 1.5) {$4(a+1/2)^2 + \frac{1}{2}(a^2-2a+1) = 1$};
\node at (2.5, 1) {$9a^2+6a+1 = 0 \qquad \therefore a=-\frac{1}{3}$};
\node at (2.5, 0.5) {$\text{から接している}$};
\end{tikzpicture}
\end{center}

\end{document}
